\documentclass[11pt, a4paper]{article}
\usepackage[a4paper, margin=2.2cm]{geometry}
\usepackage{fontspec}
\usepackage[spanish]{babel}
\usepackage{enumitem}
\setlist[itemize]{label=-}
\usepackage{amsmath}
\usepackage{amssymb}
\usepackage{booktabs}
\usepackage[table]{xcolor}
\usepackage{hyperref}

\definecolor{primary}{HTML}{1E3A8A}
\definecolor{secondary}{HTML}{475569}

\pagestyle{plain}

\begin{document}

% ==========================================
% PÁGINA 1: ENUNCIADO DEL EJERCICIO
% ==========================================
\begin{center}
    {\Large \textbf{\textcolor{primary}{FÍSICA MATEMÁTICA III}}}\\[4pt]
    {\large \textbf{Ejercicio de Certamen: Tabla de Caracteres y Relaciones de Ortogonalidad}}\\[6pt]
    {\small \textcolor{secondary}{Tiempo estimado: 25 minutos \quad | \quad Puntaje: 25 puntos}}
\end{center}
\vspace{2mm}
\hrule height 1pt
\vspace{4mm}

\subsection*{\textcolor{primary}{Enunciado del Problema}}
Considere el grupo simétrico $S_4$ (isomorfo al grupo de rotaciones propias del cubo y al grupo tetraédrico completo $T_d$), cuyo orden es $O(S_4) = 4! = 24$. Sus elementos se dividen en $s = 5$ clases de conjugación según su estructura en ciclos $[l_1^{\nu_1} \dots l_r^{\nu_r}]$:
\begin{itemize}
    \item $K_1 = [1^4]$: Contiene a la identidad $e$, con $k_1 = 1$ elemento.
    \item $K_2 = [1^2 2^1]$: Contiene las transposiciones simples $(ab)$, con $k_2 = 6$ elementos.
    \item $K_3 = [2^2]$: Contiene los productos de dos transposiciones disjuntas $(ab)(cd)$, con $k_3 = 3$ elementos.
    \item $K_4 = [1^1 3^1]$: Contiene los $3$-ciclos $(abc)$, con $k_4 = 8$ elementos.
    \item $K_5 = [4^1]$: Contiene los $4$-ciclos $(abcd)$, con $k_5 = 6$ elementos.
\end{itemize}

A continuación se muestra la tabla de caracteres parcialmente completa del grupo $S_4$, donde todos los caracteres son números reales y las incógnitas están denotadas por las letras $\{a, b, c, d, x, y, z, u, v, w\}$:

\vspace{2mm}
\begin{center}
\renewcommand{\arraystretch}{1.35}
\begin{tabular}{c|ccccc}
\toprule
$S_4$ ($O(S_4)=24$) & $K_1 = [1^4]$ & $K_2 = [1^2 2^1]$ & $K_3 = [2^2]$ & $K_4 = [1^1 3^1]$ & $K_5 = [4^1]$ \\
Cardinalidad $k_\alpha$ & $k_1 = 1$ & $k_2 = 6$ & $k_3 = 3$ & $k_4 = 8$ & $k_5 = 6$ \\
\midrule
$D^{(1)}$ & $a$ & $b$ & $c$ & $d$ & $1$ \\
$D^{(2)}$ & $1$ & $-1$ & $1$ & $1$ & $-1$ \\
$D^{(3)}$ & $x$ & $0$ & $y$ & $-1$ & $0$ \\
$D^{(4)}$ & $3$ & $1$ & $-1$ & $0$ & $z$ \\
$D^{(5)}$ & $u$ & $v$ & $w$ & $0$ & $1$ \\
\bottomrule
\end{tabular}
\end{center}
\vspace{2mm}

\noindent \textbf{Se pide:}
\begin{enumerate}[label=\textbf{\alph*)}]
    \item \textbf{(15 puntos)} Utilizando exclusivamente las propiedades generales de las tablas de caracteres (presencia de la irrep trivial, suma de los cuadrados de las dimensiones, caracteres de la representación regular y relaciones de ortogonalidad de filas y columnas), determine de manera justificada el valor numérico exacto de cada una de las 10 incógnitas $\{a, b, c, d, x, y, z, u, v, w\}$ y escriba la tabla de caracteres completa.
    \item \textbf{(5 puntos)} Verifique explícitamente que los caracteres obtenidos para las representaciones irreducibles $D^{(3)}$ y $D^{(5)}$ cumplen la condición de normalización $\langle \chi^{(3)} | \chi^{(3)} \rangle = 1$ y son mutuamente ortogonales ($\langle \chi^{(3)} | \chi^{(5)} \rangle = 0$).
    \item \textbf{(5 puntos)} Considere la representación de permutación natural $D^{\text{nat}}$ de dimensión $4$, donde cada elemento $\pi \in S_4$ permuta los 4 vectores de la base canónica $\{|e_1\rangle, |e_2\rangle, |e_3\rangle, |e_4\rangle\}$ de $\mathbb{C}^4$. Sabiendo que el carácter $\chi^{\text{nat}}(K_\alpha)$ coincide con el número de elementos que quedan fijos (ciclos de longitud $1$) en la clase $K_\alpha$, obtenga el vector de caracteres $\chi^{\text{nat}} = (4, 2, 0, 1, 0)$ y descomponga $D^{\text{nat}}$ en suma directa de representaciones irreducibles de $S_4$.
\end{enumerate}

\newpage

% ==========================================
% PÁGINA 2: SOLUCIONARIO Y REFERENCIAS AL APUNTE
% ==========================================
\begin{center}
    {\Large \textbf{\textcolor{primary}{SOLUCIONARIO Y REFERENCIAS AL APUNTE}}}
\end{center}
\vspace{1mm}
\hrule height 1pt
\vspace{3mm}

\subsection*{\textcolor{primary}{Páginas y Ecuaciones del Apunte Utilizadas}}
\begin{itemize}
    \item \textbf{Páginas 29 y 30 (Sección 2.3):} Propiedad 2 de los caracteres ($\chi^{(a)}(e) = d_a$, dimensión de la irrep) y \textbf{Teorema 66} (Teorema de ortogonalidad de caracteres, Ec. 2.56).
    \item \textbf{Páginas 31 y 32 (Sección 2.4):} Descomposición de representaciones reducibles, cálculo de multiplicidades $n_a = \langle \chi^{(a)} | \chi \rangle$ (\textbf{Ec. 2.65}), criterio de irreducibilidad $\langle \chi | \chi \rangle = 1$ (\textbf{Proposición 67, Ec. 2.66}) y propiedades de la representación regular: $\sum_{a=1}^r d_a^2 = O(G)$ (\textbf{Ec. 2.76}) y $\sum_{a=1}^r d_a \chi^{(a)}(g) = O(G)\delta_{ge}$ (\textbf{Ec. 2.77}).
    \item \textbf{Páginas 36, 37 y 38 (Sección 2.5):} \textbf{Definición 71} (Construcción de la tabla de caracteres, reglas 1 a 5):
    \begin{align*}
        \text{Suma de dimensiones (Ec. 2.114):} \quad & \sum_{a=1}^r d_a^2 = O(G) \\
        \text{Ortogonalidad por filas (Ec. 2.116):} \quad & \sum_{\alpha=1}^s k_\alpha \chi_\alpha^{*(a)} \chi_\alpha^{(b)} = O(G) \delta_{ab} \\
        \text{Ortogonalidad por columnas (Ec. 2.117):} \quad & \sum_{a=1}^r k_\alpha \chi_\alpha^{*(a)} \chi_\beta^{(a)} = O(G) \delta_{\alpha\beta} \iff \sum_{a=1}^r \chi_\alpha^{*(a)} \chi_\beta^{(a)} = \frac{O(G)}{k_\alpha} \delta_{\alpha\beta}
    \end{align*}
\end{itemize}

\subsection*{\textcolor{primary}{Desarrollo Paso a Paso}}

\noindent \textbf{Parte a): Determinación de las incógnitas}
\begin{enumerate}[label=\textbf{Paso \arabic*:}]
    \item \textbf{Representación trivial $D^{(1)}$ ($a, b, c, d$):} Por la \textbf{Regla 2 de la Definición 71 (pág. 37)}, toda tabla de caracteres posee la irrep trivial unidimensional donde todos los elementos son representados por $1$. Observando la primera fila (que termina en $1$), concluimos inmediatamente:
    \[
    a = 1, \qquad b = 1, \qquad c = 1, \qquad d = 1.
    \]
    \item \textbf{Dimensiones de las irreps $D^{(3)}$ y $D^{(5)}$ ($x, u$):} La primera columna ($K_1 = \{e\}$) contiene las dimensiones $d_a = \chi_1^{(a)}$ (\textbf{Propiedad 2, pág. 30}). Además, la columna $K_4$ tiene ceros en $D^{(4)}$ y $D^{(5)}$, por lo que al aplicar la \textbf{Regla 3 (Ec. 2.115, pág. 37)} sobre la clase $\alpha = 4$ ($\sum_{a=1}^5 d_a \chi_4^{(a)} = 0$) despejamos $x$ directamente:
    \[
    (1)(1) + (1)(1) + x(-1) + 3(0) + u(0) = 0 \implies 2 - x = 0 \implies x = 2.
    \]
    Ahora aplicamos la relación de cierre de dimensiones (\textbf{Ec. 2.76, pág. 32 / Ec. 2.114, pág. 37}), sabiendo que $u = d_5 \in \mathbb{N}$:
    \[
    \sum_{a=1}^5 d_a^2 = 1^2 + 1^2 + 2^2 + 3^2 + u^2 = 24 \implies 15 + u^2 = 24 \implies u^2 = 9 \implies u = 3.
    \]
    \item \textbf{Cálculo de $z$ en $D^{(4)}$:} Aplicamos la \textbf{ortogonalidad de filas (Ec. 2.116, pág. 38)} entre $D^{(1)}$ y $D^{(4)}$:
    \[
    \sum_{\alpha=1}^5 k_\alpha \chi_\alpha^{(1)} \chi_\alpha^{(4)} = 1(1)(3) + 6(1)(1) + 3(1)(-1) + 8(1)(0) + 6(1)(z) = 0 \implies 6 + 6z = 0 \implies z = -1.
    \]
    \item \textbf{Cálculo de $y$ en $D^{(3)}$:} Aplicamos la \textbf{ortogonalidad de filas (Ec. 2.116)} entre $D^{(1)}$ y $D^{(3)}$:
    \[
    \sum_{\alpha=1}^5 k_\alpha \chi_\alpha^{(1)} \chi_\alpha^{(3)} = 1(1)(2) + 6(1)(0) + 3(1)(y) + 8(1)(-1) + 6(1)(0) = 0 \implies 2 + 3y - 8 = 0 \implies y = 2.
    \]
    \item \textbf{Cálculo de $v$ y $w$ en $D^{(5)}$:} Aplicamos la \textbf{ortogonalidad de columnas (Ec. 2.117, pág. 38)} entre la primera columna ($\alpha = 1$, donde $\chi_1^{(a)} = d_a$) y las columnas $\beta = 2$ y $\beta = 3$ (o equivalentemente la \textbf{Ec. 2.115}):
    \begin{itemize}
        \item Columna 1 con Columna 2 ($\beta = 2$): $(1)(1) + (1)(-1) + (2)(0) + (3)(1) + 3v = 0 \implies 3 + 3v = 0 \implies v = -1$.
        \item Columna 1 con Columna 3 ($\beta = 3$): $(1)(1) + (1)(1) + (2)(2) + (3)(-1) + 3w = 0 \implies 3 + 3w = 0 \implies w = -1$.
    \end{itemize}
\end{enumerate}

\noindent \textbf{Tabla de caracteres completa de $S_4$:}
\begin{center}
\renewcommand{\arraystretch}{1.25}
\begin{tabular}{c|ccccc}
\toprule
$S_4$ & $K_1$ ($k_1=1$) & $K_2$ ($k_2=6$) & $K_3$ ($k_3=3$) & $K_4$ ($k_4=8$) & $K_5$ ($k_5=6$) \\
\midrule
$D^{(1)}$ & $1$ & $1$ & $1$ & $1$ & $1$ \\
$D^{(2)}$ & $1$ & $-1$ & $1$ & $1$ & $-1$ \\
$D^{(3)}$ & $2$ & $0$ & $2$ & $-1$ & $0$ \\
$D^{(4)}$ & $3$ & $1$ & $-1$ & $0$ & $-1$ \\
$D^{(5)}$ & $3$ & $-1$ & $-1$ & $0$ & $1$ \\
\bottomrule
\end{tabular}
\end{center}

\noindent \textbf{Parte b): Verificación de ortonormalidad entre $D^{(3)}$ y $D^{(5)}$}
Usando el producto escalar de caracteres (\textbf{Ec. 2.108, pág. 37}):
\begin{align*}
    \langle \chi^{(3)} | \chi^{(3)} \rangle &= \frac{1}{24} \sum_{\alpha=1}^5 k_\alpha \left(\chi_\alpha^{(3)}\right)^2 = \frac{1}{24} \left[ 1(2)^2 + 6(0)^2 + 3(2)^2 + 8(-1)^2 + 6(0)^2 \right] = \frac{4 + 0 + 12 + 8 + 0}{24} = \frac{24}{24} = 1. \\
    \langle \chi^{(3)} | \chi^{(5)} \rangle &= \frac{1}{24} \sum_{\alpha=1}^5 k_\alpha \chi_\alpha^{(3)} \chi_\alpha^{(5)} = \frac{1}{24} \left[ 1(2)(3) + 6(0)(-1) + 3(2)(-1) + 8(-1)(0) + 6(0)(1) \right] = \frac{6 - 6}{24} = 0.
\end{align*}

\noindent \textbf{Parte c): Reducción de $D^{\text{nat}}$ con $\chi^{\text{nat}} = (4, 2, 0, 1, 0)$}
Calculamos las multiplicidades $n_a = \langle \chi^{(a)} | \chi^{\text{nat}} \rangle$ usando la \textbf{Ecuación (2.65) (pág. 31)}:
\begin{align*}
    n_1 &= \frac{1}{24} \left[ 1(1)(4) + 6(1)(2) + 3(1)(0) + 8(1)(1) + 6(1)(0) \right] = \frac{4 + 12 + 0 + 8 + 0}{24} = 1, \\
    n_4 &= \frac{1}{24} \left[ 1(3)(4) + 6(1)(2) + 3(-1)(0) + 8(0)(1) + 6(-1)(0) \right] = \frac{12 + 12 + 0 + 0 + 0}{24} = 1,
\end{align*}
y análogamente $n_2 = n_3 = n_5 = 0$ (ya que las dimensiones ya suman $d_1 + d_4 = 1 + 3 = 4$). Por lo tanto, $D^{\text{nat}} = D^{(1)} \oplus D^{(4)}$.

\end{document}